% rope_frame_standalone.tex
% Kompiliere 20x mit: pdflatex "\def\pos{0}\input{rope_frame_standalone.tex}"
%                      pdflatex "\def\pos{1}\input{rope_frame_standalone.tex}"
%                      ... etc.
% ODER nutze das Shell-Skript unten

\ifdefined\pos\else\def\pos{0}\fi

\documentclass[border=5pt]{standalone}
\usepackage[utf8]{inputenc}
\usepackage[T1]{fontenc}
\usepackage{tikz}
\usetikzlibrary{arrows.meta}

\begin{document}

\definecolor{dim1}{HTML}{E8590C}
\definecolor{dim2}{HTML}{0F3460}
\definecolor{dim3}{HTML}{E94560}
\definecolor{dim4}{HTML}{1A8F3C}
\definecolor{dim5}{HTML}{7B2D8B}
\definecolor{dim6}{HTML}{D4A017}

\begin{tikzpicture}
  \useasboundingbox (-6.5,-6.5) rectangle (6.5,6.5);
  \fill[white] (-6.5,-6.5) rectangle (6.5,6.5);

  % ── Winkel berechnen ──
  \pgfmathsetmacro{\aA}{36*\pos}
  \pgfmathsetmacro{\aB}{24*\pos}
  \pgfmathsetmacro{\aC}{16*\pos}
  \pgfmathsetmacro{\aD}{10*\pos}
  \pgfmathsetmacro{\aE}{6*\pos}
  \pgfmathsetmacro{\aF}{3*\pos}

  % ── Referenzlinien ──
  \draw[gray!15, thin] (0,0) -- (0:5.8);
  \draw[gray!15, thin] (0,0) -- (90:5.8);
  \draw[gray!15, thin] (0,0) -- (180:5.8);
  \draw[gray!15, thin] (0,0) -- (270:5.8);

  % ── Kreise ──
  \draw[dim1, thick, opacity=0.5] (0,0) circle (1.3);
  \draw[dim2, thick, opacity=0.45] (0,0) circle (2.1);
  \draw[dim3, thick, opacity=0.4] (0,0) circle (2.9);
  \draw[dim4, thick, opacity=0.35] (0,0) circle (3.7);
  \draw[dim5, thick, opacity=0.3] (0,0) circle (4.5);
  \draw[dim6, thick, opacity=0.25] (0,0) circle (5.3);

  % ── Labels ──
  \node[font=\tiny, dim1] at (200:1.6) {$(d_0,d_1)$};
  \node[font=\tiny, dim2] at (200:2.4) {$(d_2,d_3)$};
  \node[font=\tiny, dim3] at (200:3.2) {$(d_4,d_5)$};
  \node[font=\tiny, dim4] at (200:4.0) {$(d_6,d_7)$};
  \node[font=\tiny, dim5] at (200:4.8) {$(d_8,d_9)$};
  \node[font=\tiny, dim6] at (200:5.6) {$(d_{10},d_{11})$};

  % ── Verbindungslinien ──
  \draw[gray, thin, opacity=0.15] (0,0) -- (\aA:1.3);
  \draw[gray, thin, opacity=0.12] (0,0) -- (\aB:2.1);
  \draw[gray, thin, opacity=0.1] (0,0) -- (\aC:2.9);
  \draw[gray, thin, opacity=0.08] (0,0) -- (\aD:3.7);
  \draw[gray, thin, opacity=0.06] (0,0) -- (\aE:4.5);
  \draw[gray, thin, opacity=0.04] (0,0) -- (\aF:5.3);

  % ── Spur (gesamter Bogen von 0 bis aktuelle Position) ──
  \ifnum\pos>0\relax
    \draw[dim1, thick, opacity=0.2] (0:1.3) arc (0:\aA:1.3);
    \draw[dim2, thick, opacity=0.2] (0:2.1) arc (0:\aB:2.1);
    \draw[dim3, thick, opacity=0.2] (0:2.9) arc (0:\aC:2.9);
    \draw[dim4, thick, opacity=0.2] (0:3.7) arc (0:\aD:3.7);
    \draw[dim5, thick, opacity=0.2] (0:4.5) arc (0:\aE:4.5);
    \draw[dim6, thick, opacity=0.2] (0:5.3) arc (0:\aF:5.3);
  \fi

  % ── Glow ──
  \fill[dim1, opacity=0.12] (\aA:1.3) circle (9pt);
  \fill[dim2, opacity=0.12] (\aB:2.1) circle (9pt);
  \fill[dim3, opacity=0.12] (\aC:2.9) circle (9pt);
  \fill[dim4, opacity=0.12] (\aD:3.7) circle (9pt);
  \fill[dim5, opacity=0.12] (\aE:4.5) circle (9pt);
  \fill[dim6, opacity=0.12] (\aF:5.3) circle (9pt);

  % ── Punkte ──
  \fill[dim1] (\aA:1.3) circle (5pt);
  \fill[dim2] (\aB:2.1) circle (5pt);
  \fill[dim3] (\aC:2.9) circle (5pt);
  \fill[dim4] (\aD:3.7) circle (5pt);
  \fill[dim5] (\aE:4.5) circle (5pt);
  \fill[dim6] (\aF:5.3) circle (5pt);

  % ── Startpunkte (pos=0) ──
  \draw[gray!40] (0:1.3) circle (2pt);
  \draw[gray!40] (0:2.1) circle (2pt);
  \draw[gray!40] (0:2.9) circle (2pt);
  \draw[gray!40] (0:3.7) circle (2pt);
  \draw[gray!40] (0:4.5) circle (2pt);
  \draw[gray!40] (0:5.3) circle (2pt);

  % ── Wort ──
  \fill[white] (0,0) circle (0.75);
  \node[draw, thick, rounded corners=4pt, fill=black!85, text=white,
        minimum width=1.4cm, minimum height=0.6cm,
        font=\bfseries] at (0,0) {Katze};

  % ── Position-Anzeige ──
  \node[font=\Large\bfseries, anchor=north west] at (-6.2, 6.2)
    {pos = \pos};

  % ── Legende ──
  \node[anchor=north east, font=\tiny, text width=2.8cm, align=left,
        fill=white, draw=gray!30, rounded corners, inner sep=3pt]
    at (6.3, 6.2) {%
    \textcolor{dim1}{$\bullet$ 36\textdegree/pos (schnell)}\\
    \textcolor{dim2}{$\bullet$ 24\textdegree/pos}\\
    \textcolor{dim3}{$\bullet$ 16\textdegree/pos}\\
    \textcolor{dim4}{$\bullet$ 10\textdegree/pos}\\
    \textcolor{dim5}{$\bullet$ 6\textdegree/pos}\\
    \textcolor{dim6}{$\bullet$ 3\textdegree/pos (langsam)}
  };

  % ── Untertitel ──
  \node[anchor=south, font=\scriptsize, text width=11cm, align=center]
    at (0, -6.3) {%
    Jeder Punkt rotiert mit eigener Frequenz.
    Zusammen = eindeutiger Positions-Fingerabdruck.
  };

\end{tikzpicture}
\end{document}
